\documentclass[11pt]{article}

\usepackage[margin=1in]{geometry}
\usepackage{amsmath, amssymb}
\usepackage{booktabs}
\usepackage{array}
\usepackage{graphicx}
\usepackage{xcolor}
\usepackage{hyperref}
\usepackage{enumitem}
\usepackage{tcolorbox}
\usepackage{caption}

\hypersetup{colorlinks=true, linkcolor=blue, urlcolor=blue, citecolor=blue}

\setlist[itemize]{noitemsep, topsep=2pt}
\setlist[enumerate]{noitemsep, topsep=2pt}

\newcommand{\bhat}{\widehat{\beta}}
\newcommand{\yhat}{\widehat{Y}}
\newcommand{\uhat}{\widehat{u}}

\title{\textbf{OLS Regression in Excel}\\[0.9em]
\large ECON 300 - Labour Economics}
\author{}
\date{}

\begin{document}
\maketitle

\section*{Learning Objectives}
By the end of these notes, you will be able to:
\begin{itemize}
  \item Explain what a regression is and what OLS (Ordinary Least Squares) does.
  \item Set up consumption and income data in Excel.
  \item Make a scatter plot and add a regression (trendline) in Excel.
  \item Run OLS using Excel's \emph{Regression} tool (Analysis ToolPak).
  \item Interpret key outputs: coefficients, $R^2$, and p-values.
  \item Compute predicted values and residuals and create a residual plot.
\end{itemize}

\section*{What is Regression? (Plain Language)}
A \textbf{regression} is a method used to summarize how an outcome variable changes with a predictor variable.

\begin{tcolorbox}[colback=gray!5, colframe=black!40, title=Key Idea]
Regression fits a line through a scatter plot of data points to describe the relationship between two variables.
\end{tcolorbox}

In this note, we study how \textbf{Consumption} depends on \textbf{Income}.

\section*{Variables and the Model}
\subsection*{Definitions}
\begin{itemize}
  \item \textbf{Income} ($Y_i$): the income for observation $i$ (e.g., household $i$).
  \item \textbf{Consumption} ($C_i$): the consumption for observation $i$.
  \item \textbf{Observation} ($i$): one unit in the data (e.g., one household).
\end{itemize}

\subsection*{The Simple Linear Regression Model}
We write:
\begin{equation}
C_i = \beta_0 + \beta_1 Y_i + u_i,
\end{equation}
where:
\begin{itemize}
  \item $\beta_0$ is the \textbf{intercept} (also called \textbf{autonomous consumption}).
  \item $\beta_1$ is the \textbf{slope} (interpreted as the \textbf{marginal propensity to consume}, MPC).
  \item $u_i$ is the \textbf{error term}: everything affecting consumption besides income.
\end{itemize}

\subsection*{Interpretation}
\begin{itemize}
  \item If $\beta_1 > 0$, higher income is associated with higher consumption.
  \item The \textbf{MPC} ($\beta_1$) is: ``If income increases by 1 unit, consumption increases by $\beta_1$ units (on average).'' 
\end{itemize}

\section*{What Does OLS Do?}
\textbf{OLS} (Ordinary Least Squares) chooses $\bhat_0$ and $\bhat_1$ to make the fitted line as close as possible to the data points.

\subsection*{Fitted Values and Residuals}
Given estimates $\bhat_0$ and $\bhat_1$, the predicted consumption is:
\begin{equation}
\widehat{C}_i = \bhat_0 + \bhat_1 Y_i.
\end{equation}
The \textbf{residual} is:
\begin{equation}
\uhat_i = C_i - \widehat{C}_i.
\end{equation}
Residuals are the vertical gaps between the points and the fitted line.

\subsection*{Least Squares Criterion}
OLS picks $\bhat_0$ and $\bhat_1$ to minimize the \textbf{Sum of Squared Residuals (SSR)}:
\begin{equation}
\text{SSR} = \sum_{i=1}^n \uhat_i^2 = \sum_{i=1}^n (C_i - \bhat_0 - \bhat_1 Y_i)^2.
\end{equation}

\section*{Step 1: Enter Consumption and Income Data in Excel}
Open Excel and create two columns:

\begin{center}
\begin{tabular}{>{\raggedleft\arraybackslash}p{0.22\linewidth} >{\raggedleft\arraybackslash}p{0.22\linewidth}}
\toprule
\textbf{Income} ($Y$) & \textbf{Consumption} ($C$) \\
\midrule
20 & 18 \\
25 & 21 \\
30 & 24 \\
35 & 27 \\
40 & 30 \\
45 & 33 \\
50 & 36 \\
55 & 38 \\
\bottomrule
\end{tabular}
\end{center}

\begin{tcolorbox}[colback=gray!5, colframe=black!40, title=Excel Setup]
\begin{enumerate}
  \item In cell \texttt{A1} type \texttt{Income}.
  \item In cell \texttt{B1} type \texttt{Consumption}.
  \item Enter the values in rows 2--9.
\end{enumerate}
\end{tcolorbox}

\paragraph{Units.} You can interpret the numbers as ``thousand dollars'' to keep the scale simple.

\section*{Step 2: Make a Scatter Plot (Graphical Illustration)}
A scatter plot shows the relationship visually.

\begin{tcolorbox}[colback=gray!5, colframe=black!40, title=Why we plot first]
Before running regression, always plot the data to see whether a linear relationship is reasonable.
\end{tcolorbox}

\subsection*{Excel Steps (Scatter Plot)}
\begin{enumerate}
  \item Highlight the range \texttt{A1:B9}.
  \item Go to \textbf{Insert} $\rightarrow$ \textbf{Charts} $\rightarrow$ \textbf{Scatter} $\rightarrow$ \textbf{Scatter with only markers}.
  \item Ensure the horizontal axis is \textbf{Income} and the vertical axis is \textbf{Consumption}.
\end{enumerate}

\subsection*{Add Axis Titles (Recommended)}
\begin{enumerate}
  \item Click the chart.
  \item Use \textbf{Chart Elements} (the ``+'' icon in Windows) or the chart formatting menu (Mac).
  \item Add \textbf{Axis Titles}:
  \begin{itemize}
    \item Horizontal: \texttt{Income}
    \item Vertical: \texttt{Consumption}
  \end{itemize}
\end{enumerate}

\section*{Step 3: Add a Linear Trendline (Quick Regression Line)}
Excel can draw the OLS line directly on the plot using a \textbf{trendline}.

\subsection*{Excel Steps (Trendline + Equation + $R^2$)}
\begin{enumerate}
  \item Click one of the data points in the scatter plot.
  \item Right-click $\rightarrow$ \textbf{Add Trendline}.
  \item Choose \textbf{Linear}.
  \item Check:
  \begin{itemize}
    \item \textbf{Display Equation on chart}
    \item \textbf{Display R-squared value on chart}
  \end{itemize}
\end{enumerate}

\subsection*{What the equation means}
If Excel shows something like:
\[
C = 6.5 + 0.583Y,
\]
then:
\begin{itemize}
  \item Intercept $6.5$ = predicted consumption when income is $0$.
  \item Slope $0.583$ = MPC: when income rises by 1 unit, predicted consumption rises by $0.583$ units.
\end{itemize}

\section*{Step 4: Enable the Analysis ToolPak (Regression Output Table)}
To get a full regression table (coefficients, standard errors, p-values), you need the \textbf{Analysis ToolPak}.

\subsection*{Windows}
\begin{enumerate}
  \item \textbf{File} $\rightarrow$ \textbf{Options} $\rightarrow$ \textbf{Add-ins}.
  \item At the bottom: \textbf{Manage: Excel Add-ins} $\rightarrow$ \textbf{Go}.
  \item Check \textbf{Analysis ToolPak} $\rightarrow$ \textbf{OK}.
\end{enumerate}

\subsection*{Mac}
\begin{enumerate}
  \item \textbf{Tools} $\rightarrow$ \textbf{Excel Add-ins}.
  \item Check \textbf{Analysis ToolPak} $\rightarrow$ \textbf{OK}.
\end{enumerate}

\section*{Step 5: Run OLS Using Data Analysis $\rightarrow$ Regression}
\subsection*{Excel Steps (Regression Tool)}
\begin{enumerate}
  \item Go to \textbf{Data} tab.
  \item Click \textbf{Data Analysis}.
  \item Select \textbf{Regression} $\rightarrow$ \textbf{OK}.
\end{enumerate}

You will see a dialog with boxes. Fill them in as follows:

\begin{tcolorbox}[colback=gray!5, colframe=black!40, title=Regression Dialog Settings]
\begin{itemize}
  \item \textbf{Input Y Range} (Outcome): \texttt{B1:B9} \quad (Consumption)
  \item \textbf{Input X Range} (Predictor): \texttt{A1:A9} \quad (Income)
  \item Check \textbf{Labels} (because row 1 has headers).
  \item Choose \textbf{Output Range}: e.g., \texttt{D1} (or select ``New Worksheet Ply'').
  \item Optional (recommended): check \textbf{Residuals}.
\end{itemize}
\end{tcolorbox}

\section*{Step 6: Understand the Regression Output}
Excel produces several blocks. Focus on the following.

\subsection*{A) Coefficients Table (Most Important)}
You will see rows like:
\begin{itemize}
  \item \textbf{Intercept}
  \item \textbf{Income}
\end{itemize}
and columns like:
\begin{itemize}
  \item \textbf{Coefficient}
  \item \textbf{Standard Error}
  \item \textbf{t Stat}
  \item \textbf{P-value}
\end{itemize}

\paragraph{Coefficient (Interpretation).}
\begin{itemize}
  \item Intercept coefficient $\bhat_0$: predicted consumption when income is 0.
  \item Income coefficient $\bhat_1$: predicted change in consumption for a 1-unit increase in income (the MPC).
\end{itemize}

\paragraph{Standard Error (Uncertainty).}
A standard error measures how uncertain the coefficient estimate is. Smaller standard error means more precise estimate.

\paragraph{t-statistic.}
\[
t = \frac{\text{Coefficient}}{\text{Standard Error}}.
\]
Bigger $|t|$ means stronger evidence the coefficient is different from zero.

\paragraph{p-value.}
The p-value answers:
\begin{quote}
If the true effect were 0, how likely would we see a coefficient this large (in absolute value) just by chance?
\end{quote}
Rule of thumb:
\[
\text{If p-value } < 0.05, \text{ we say the coefficient is statistically significant.}
\]

\subsection*{B) $R^2$ (Goodness of Fit)}
$R^2$ is the fraction of variation in consumption explained by income (in a linear way).
\[
0 \le R^2 \le 1.
\]
Interpretation examples:
\begin{itemize}
  \item $R^2 = 0.80$: income explains 80\% of the variation in consumption.
  \item $R^2 = 0.10$: income explains only 10\% (weak linear fit).
\end{itemize}

\subsection*{C) ``Significance F'' (Overall Test)}
In a one-predictor model, this is closely related to testing whether the slope is zero.
A small value (e.g., $<0.05$) suggests the regression model explains consumption better than a model with no predictor.

\newpage
\section*{Step 7: Compute Predicted Consumption and Residuals in Excel}
Even if Excel outputs residuals, you should know how to compute them.

\subsection*{A) Predicted Consumption}
Suppose your regression output shows:
\[
\bhat_0 = \text{(Intercept)}, \quad \bhat_1 = \text{(Income coefficient)}.
\]
Store them in two cells. For example:
\begin{itemize}
  \item Put $\bhat_0$ in cell \texttt{E2}
  \item Put $\bhat_1$ in cell \texttt{E3}
\end{itemize}
(Or just reference the cells where Excel printed them.)

In column \texttt{C}:
\begin{itemize}
  \item \texttt{C1}: \texttt{Predicted\_C}
  \item \texttt{C2}: type the formula:
\end{itemize}

\begin{tcolorbox}[colback=gray!5, colframe=black!40, title=Excel Formula: Predicted Consumption]
\texttt{= \$E\$2 + \$E\$3 * A2}
\end{tcolorbox}

Copy the formula down to \texttt{C9}.

\subsection*{B) Residuals}
In column \texttt{D}:
\begin{itemize}
  \item \texttt{D1}: \texttt{Residual}
  \item \texttt{D2}: type:
\end{itemize}

\begin{tcolorbox}[colback=gray!5, colframe=black!40, title=Excel Formula: Residual]
\texttt{= B2 - C2}
\end{tcolorbox}

Copy the formula down to \texttt{D9}.

\section*{Step 8: Make the Residual Plot (Model Check)}
A residual plot checks whether a straight line is a reasonable summary.

\subsection*{Excel Steps (Residual Plot)}
\begin{enumerate}
  \item Highlight \texttt{A1:A9} (Income).
  \item Hold Ctrl (Windows) or Cmd (Mac) and also select \texttt{D1:D9} (Residuals).
  \item Insert $\rightarrow$ Scatter (only markers).
  \item Title the chart: \texttt{Residuals vs Income}.
\end{enumerate}

\subsection*{How to interpret the residual plot}
\begin{itemize}
  \item \textbf{Good sign:} residuals look randomly scattered around 0.
  \item \textbf{Warning sign:} a curve pattern (suggests nonlinearity).
  \item \textbf{Warning sign:} ``fan shape'' (spread increasing with income) suggests unequal variance.
\end{itemize}

\section*{Alternative Method: Regression Using Excel Functions (No ToolPak)}
Excel can compute slope, intercept, and $R^2$ using built-in functions.

\subsection*{Slope (MPC)}
\begin{tcolorbox}[colback=gray!5, colframe=black!40, title=Excel Function: SLOPE]
\texttt{=SLOPE(B2:B9, A2:A9)}
\end{tcolorbox}

\subsection*{Intercept (Autonomous Consumption)}
\begin{tcolorbox}[colback=gray!5, colframe=black!40, title=Excel Function: INTERCEPT]
\texttt{=INTERCEPT(B2:B9, A2:A9)}
\end{tcolorbox}

\subsection*{$R^2$}
\begin{tcolorbox}[colback=gray!5, colframe=black!40, title=Excel Function: RSQ]
\texttt{=RSQ(B2:B9, A2:A9)}
\end{tcolorbox}

\section*{How to Write the Final Answer (Economic Interpretation)}
Once you have estimates, write a sentence like:

\begin{tcolorbox}[colback=gray!5, colframe=black!40, title=Interpretation Template]
The estimated consumption function is $\widehat{C} = \bhat_0 + \bhat_1 Y$. The slope $\bhat_1$ is the MPC: a one-unit increase in income is associated with an average increase of $\bhat_1$ units in consumption. The intercept $\bhat_0$ is autonomous consumption (predicted consumption when income is zero). The $R^2$ indicates how much of the variation in consumption is explained by income.
\end{tcolorbox}

\section*{Summary}
\begin{itemize}
  \item Regression fits a line: $C = \beta_0 + \beta_1 Y + u$.
  \item OLS chooses coefficients to minimize squared residuals.
  \item $\bhat_1$ (slope) is the MPC.
  \item Always: (1) scatter plot, (2) regression output, (3) residual plot.
\end{itemize}

\end{document}
